% ============================================================
%  CUADERNILLO DE ESTUDIO — CLASE 4
%  Seguridad de la Información y de Sistemas
%  Lic. en Gestión de Negocios Digitales — UCA
%  Compilar con: pdflatex (dos pasadas)
% ============================================================

\documentclass[12pt,a4paper]{article}

% ---------- Paquetes ----------
\usepackage[utf8]{inputenc}
\usepackage[spanish]{babel}
\usepackage[T1]{fontenc}
\usepackage{lmodern}
\usepackage[margin=2.2cm]{geometry}
\usepackage{amssymb}
\usepackage{enumitem}
\usepackage{booktabs}
\usepackage{array}
\usepackage{tabularx}
\usepackage{xcolor}
\usepackage{tcolorbox}
\usepackage{graphicx}
\usepackage{fancyhdr}
\usepackage{titlesec}
\usepackage{hyperref}
\usepackage{multicol}
\usepackage{pifont}
\usepackage{wasysym}
\usepackage{tikz}
\usetikzlibrary{positioning,arrows.meta,shapes.geometric}

% ---------- Colores ----------
\definecolor{azulUCA}{HTML}{003366}
\definecolor{doradoUCA}{HTML}{B8860B}
\definecolor{grisClaro}{HTML}{F2F2F2}
\definecolor{rojoAlerta}{HTML}{C0392B}
\definecolor{verdeOK}{HTML}{27AE60}
\definecolor{azulCielo}{HTML}{2980B9}
\definecolor{violetaIA}{HTML}{8E44AD}
\definecolor{naranjaAmenaza}{HTML}{D35400}

% ---------- tcolorbox estilos ----------
\tcbuselibrary{skins,breakable}

\newtcolorbox{cajaTitulo}[1][]{
  colback=azulUCA, colframe=azulUCA, coltext=white,
  fonttitle=\Large\bfseries, arc=4mm, #1
}

\newtcolorbox{cajaDefinicion}{
  colback=azulCielo!10, colframe=azulCielo,
  fonttitle=\bfseries, title=Definici\'on, arc=3mm, breakable
}

\newtcolorbox{cajaConcepto}{
  colback=azulUCA!8, colframe=azulUCA,
  fonttitle=\bfseries\color{azulUCA}, arc=3mm, breakable
}

\newtcolorbox{cajaEjemplo}{
  colback=verdeOK!10, colframe=verdeOK,
  fonttitle=\bfseries\color{verdeOK}, title=Ejemplo, arc=3mm, breakable
}

\newtcolorbox{cajaActividad}{
  colback=grisClaro, colframe=azulUCA,
  fonttitle=\bfseries\color{azulUCA},
  title={\ding{45}\; Actividad Pr\'actica}, arc=3mm, breakable
}

\newtcolorbox{cajaIA}{
  colback=violetaIA!8, colframe=violetaIA,
  fonttitle=\bfseries\color{violetaIA},
  title={\ding{211}\; Actividad con Inteligencia Artificial}, arc=3mm, breakable
}

\newtcolorbox{cajaNota}{
  colback=rojoAlerta!8, colframe=rojoAlerta,
  fonttitle=\bfseries\color{rojoAlerta}, title=Importante, arc=3mm, breakable
}

\newtcolorbox{cajaCierre}{
  colback=doradoUCA!10, colframe=doradoUCA,
  fonttitle=\bfseries, title=Cierre de Unidad, arc=3mm, breakable
}

% ---------- Header / Footer ----------
\pagestyle{fancy}
\fancyhf{}
\fancyhead[L]{\small\textcolor{azulUCA}{Seguridad de la Informaci\'on y de Sistemas}}
\fancyhead[R]{\small\textcolor{azulUCA}{Cuadernillo --- Clase 4}}
\fancyfoot[C]{\thepage}
\fancyfoot[L]{\small\textcolor{gray}{UCA}}
\renewcommand{\headrulewidth}{0.4pt}
\renewcommand{\footrulewidth}{0.4pt}

% ---------- Títulos ----------
\titleformat{\section}
  {\color{azulUCA}\Large\bfseries}
  {\thesection.}{0.6em}{}
\titleformat{\subsection}
  {\color{azulCielo}\large\bfseries}
  {\thesubsection}{0.6em}{}

% ---------- Enlaces ----------
\hypersetup{colorlinks=true, linkcolor=azulUCA, urlcolor=azulCielo}

% ---------- Checkbox ----------
\newcommand{\checkboxVacio}{$\square$\;}

% ============================================================
\begin{document}

% ============================================================
%  PORTADA
% ============================================================
\begin{titlepage}
\centering
\vspace*{1cm}

\begin{cajaTitulo}[width=\textwidth, halign=center]
  SEGURIDAD DE LA INFORMACI\'ON Y DE SISTEMAS\\[4pt]
  {\large Licenciatura en Gesti\'on de Negocios Digitales --- UCA}
\end{cajaTitulo}

\vspace{1.5cm}

{\Huge\bfseries\color{azulUCA} CUADERNILLO DE ESTUDIO}\\[8pt]
{\LARGE\color{doradoUCA} CLASE 4}\\[12pt]
{\Large\itshape ``Criptografía aplicada\\[2pt] a negocios digitales''}\\[0.8cm]
{\large\color{violetaIA}\ding{211}\; Incluye actividad con Inteligencia Artificial}\\[1.5cm]

\begin{tabular}{rl}
\textbf{Profesor:}   & Mg.\ Ing.\ Rodrigo Atilio Elgueta \\[4pt]
\textbf{Unidad:}     & Cierre de Unidad 1 \\[4pt]
\end{tabular}

\vfill

\begin{tcolorbox}[colback=grisClaro, colframe=azulUCA, arc=3mm, width=0.85\textwidth]
  \centering
  {\bfseries Datos del/la estudiante}\\[8pt]
  \begin{tabular}{p{5cm}l p{5cm}l}
    Nombre y apellido: & \underline{\hspace{4.8cm}} &
    Grupo N.\textdegree: & \underline{\hspace{4.8cm}} \\[10pt]
    Correo: & \underline{\hspace{4.8cm}} &
    Fecha: & \underline{\hspace{4.8cm}} \\
  \end{tabular}
\end{tcolorbox}

\vspace{1cm}
{\small\textcolor{gray}{Material de uso exclusivo para la cursada.}}
\end{titlepage}

\tableofcontents
\newpage

% ============================================================
%  SECCIÓN 1 — INTRODUCCIÓN
% ============================================================
\section{¿Por qu\'e la criptograf\'ia sostiene la confianza digital?}

En las clases anteriores vimos \textbf{qu\'e} proteger (activos, pilares CID) y \textbf{contra qu\'e} (amenazas). Hoy vamos a ver \textbf{c\'omo} se protege la informaci\'on en tr\'ansito y en reposo: la \textbf{criptograf\'ia}.

\begin{cajaConcepto}[title={Conexi\'on con lo anterior}]
Recordemos la demo de la Clase 1: la página accedi\'o a tu cámara y transmiti\'o la imagen a un panel de administraci\'on. Si esa transmisi\'on hubiera estado \textbf{cifrada}, un tercero no habr\'ia podido interceptarla. Si el panel hubiera requerido \textbf{autenticaci\'on fuerte}, no cualquiera habr\'ia podido verla. La criptograf\'ia es la base t\'ecnica que hace posibles esos controles.
\end{cajaConcepto}

\begin{cajaNota}
\textbf{Sin criptograf\'ia no hay e-commerce, no hay fintech, no hay home banking.} Cada vez que un cliente ingresa los datos de su tarjeta en tu tienda online, conf\'ia en que esa informaci\'on viaja protegida. Esa protecci\'on es criptograf\'ia.
\end{cajaNota}

\vspace{6pt}
\noindent\textbf{Objetivos de esta clase:}
\begin{enumerate}[leftmargin=2em]
  \item Comprender qu\'e es \textbf{cifrado} y \textbf{hashing}, y por qu\'e sostienen la confianza digital (HTTPS, contrase\~nas, firma digital).
  \item Practicar el uso \textbf{cr\'itico} de IA generativa para explicar conceptos t\'ecnicos.
\end{enumerate}

\newpage
% ============================================================
%  SECCIÓN 2 — CIFRADO SIMÉTRICO VS ASIMÉTRICO
% ============================================================
\section{Cifrado sim\'etrico y asim\'etrico}

\begin{cajaDefinicion}
\textbf{Cifrar} es transformar informaci\'on legible (texto plano) en una forma ilegible (texto cifrado) usando una \textbf{clave}, de modo que solo quien posee la clave adecuada pueda revertir el proceso.
\end{cajaDefinicion}

% --- 2.1 Simétrico ---
\subsection{Cifrado sim\'etrico: una sola clave}

\begin{cajaConcepto}[title={Analog\'ia del candado}]
Imagin\'a una caja fuerte con \textbf{una sola llave}. El que guarda el documento usa esa llave para cerrar la caja. El que necesita leer el documento debe tener \textbf{la misma llave} para abrirla.\\[4pt]
\textbf{Problema:} ¿C\'omo le hago llegar la llave al destinatario sin que alguien la copie en el camino?
\end{cajaConcepto}

\textbf{Caracter\'isticas:}
\begin{itemize}[leftmargin=2em]
  \item Una \'unica clave para cifrar y descifrar.
  \item Muy r\'apido y eficiente para grandes vol\'umenes de datos.
  \item El desaf\'io es la \textbf{distribuci\'on segura de la clave}.
\end{itemize}

\begin{cajaEjemplo}
\textbf{En un negocio digital:} Cifrar la base de datos de clientes en el servidor. El servidor cifra al guardar y descifra al leer, usando la misma clave que est\'a almacenada de forma segura (idealmente en un m\'odulo HSM o un gestor de secretos como AWS Secrets Manager).
\end{cajaEjemplo}

% --- 2.2 Asimétrico ---
\subsection{Cifrado asim\'etrico: dos claves (p\'ublica y privada)}

\begin{cajaConcepto}[title={Analog\'ia del buz\'on}]
Imagin\'a un buz\'on p\'ublico en la calle. \textbf{Cualquiera} puede meter una carta por la ranura (usar la \textbf{clave p\'ublica} para cifrar), pero \textbf{solo el due\~no del buz\'on} tiene la llave para abrirlo y sacar las cartas (usar la \textbf{clave privada} para descifrar).\\[4pt]
\textbf{Ventaja:} No necesito que el remitente tenga una llave secreta; solo necesita saber que el buz\'on existe.
\end{cajaConcepto}

\textbf{Caracter\'isticas:}
\begin{itemize}[leftmargin=2em]
  \item Par de claves: una \textbf{p\'ublica} (se comparte) y una \textbf{privada} (se guarda en secreto).
  \item Lo que cifra una, solo lo descifra la otra.
  \item M\'as lento que el sim\'etrico (se usa para establecer la conexi\'on, no para cifrar todo el tr\'afico).
\end{itemize}

\begin{cajaEjemplo}
\textbf{En un negocio digital:} Cuando un cliente se conecta a tu e-commerce por HTTPS, su navegador usa la \textbf{clave p\'ublica} de tu servidor para negociar una clave de sesi\'on sim\'etrica. A partir de ah\'i, toda la comunicaci\'on viaja cifrada con esa clave sim\'etrica (r\'apida y segura).
\end{cajaEjemplo}

% --- Tabla comparativa ---
\subsection{Comparaci\'on r\'apida}

\renewcommand{\arraystretch}{1.5}
\begin{tabularx}{\textwidth}{|>{\bfseries}p{3cm}|X|X|}
\hline
 & \textbf{Sim\'etrico} & \textbf{Asim\'etrico} \\
\hline
Claves & Una sola (compartida) & Par: p\'ublica + privada \\
\hline
Velocidad & Muy r\'apido & M\'as lento \\
\hline
Uso típico & Cifrar datos en reposo, sesiones HTTPS & Establecer conexi\'on, firma digital \\
\hline
Desaf\'io & ¿C\'omo comparto la clave de forma segura? & ¿C\'omo s\'e que la clave p\'ublica es realmente de quien dice ser? (certificados) \\
\hline
Ejemplo real & AES-256 (base de datos, discos) & RSA, ECC (HTTPS, firma digital) \\
\hline
\end{tabularx}

\begin{cajaNota}
\textbf{En la pr\'actica se usan juntos:} HTTPS usa cifrado asim\'etrico para ``presentarse'' y negociar, y luego sim\'etrico para transmitir los datos r\'apido. Es un sistema h\'ibrido.
\end{cajaNota}

\newpage
% ============================================================
%  SECCIÓN 3 — HASHING
% ============================================================
\section{Hashing: la huella digital de los datos}

\begin{cajaDefinicion}
Una \textbf{función hash} toma cualquier dato (una contrase\~na, un archivo, un documento) y produce una cadena de caracteres de longitud fija, \'unica para ese dato. Es \textbf{irreversible}: no se puede ``descifrar'' un hash para obtener el dato original.
\end{cajaDefinicion}

\subsection{¿Por qu\'e no es lo mismo que cifrar?}

\renewcommand{\arraystretch}{1.5}
\begin{tabularx}{\textwidth}{|>{\bfseries}p{2.5cm}|X|X|}
\hline
 & \textbf{Cifrado} & \textbf{Hash} \\
\hline
¿Reversible? & S\'i (con la clave correcta) & No (irreversible) \\
\hline
Prop\'osito & Proteger datos que necesito leer despu\'es & Verificar integridad; almacenar contrase\~nas \\
\hline
Salida & Texto cifrado de longitud variable & Cadena de longitud fija (ej. 64 chars en SHA-256) \\
\hline
Ejemplo & Cifrar un contrato PDF & Guardar la contrase\~na del usuario en la base de datos \\
\hline
\end{tabularx}

\subsection{SHA-256 en la pr\'actica}

\begin{cajaEjemplo}
Si aplico SHA-256 a la palabra \texttt{Hola}:
\begin{verbatim}
Hola -> 7426db3a45...
\end{verbatim}
Si aplico SHA-256 a \texttt{hola} (min\'uscula):
\begin{verbatim}
hola -> 11254a2fc8...
\end{verbatim}
\textbf{Un solo bit de diferencia} produce un hash completamente distinto. Eso lo hace ideal para detectar si un archivo fue modificado.
\end{cajaEjemplo}

\subsection{¿Por qu\'e nunca almacenar contrase\~nas en texto plano?}

\begin{cajaNota}
Si tu e-commerce guarda las contrase\~nas de los usuarios en texto plano y un atacante roba la base de datos, \textbf{obtiene todas las contrase\~nas directamente}.\\[4pt]
Si en cambio guard\'as el \textbf{hash} de cada contrase\~na, el atacante obtiene cadenas ilegibles. Para ``adivinar'' la contraseña original tendr\'ia que probar millones de combinaciones una por una (ataque de fuerza bruta o diccionario).\\[4pt]
\textbf{Regla:} Las contrase\~nas se almacenan hasheadas (idealmente con \emph{salt} y funciones lentas como bcrypt o Argon2). Nunca en texto plano. Nunca con hash simple como MD5.
\end{cajaNota}

\newpage
% ============================================================
%  SECCIÓN 4 — HTTPS, CERTIFICADOS Y FIRMA DIGITAL
% ============================================================
\section{HTTPS, certificados y firma digital}

% --- HTTPS ---
\subsection{HTTPS: el candado del navegador}

\begin{cajaConcepto}[title={¿Qu\'e pasa cuando ves el candadito?}]
Cuando un cliente entra a \texttt{https://mitienda.com}, ocurre esto:
\begin{enumerate}[leftmargin=2em]
  \item El navegador pregunta al servidor: ``¿Qui\'en sos?''.
  \item El servidor responde con su \textbf{certificado digital} (emitido por una Autoridad Certificante como Let's Encrypt, DigiCert, etc.).
  \item El navegador verifica que el certificado es v\'alido y confiable.
  \item Se negocia una \textbf{clave de sesi\'on sim\'etrica} usando criptograf\'ia asim\'etrica.
  \item Toda la comunicaci\'on posterior viaja \textbf{cifrada} con esa clave.
\end{enumerate}
\textbf{Resultado:} Nadie en el medio (ni en el WiFi del café) puede leer los datos de la tarjeta del cliente.
\end{cajaConcepto}

\begin{cajaNota}
\textbf{Para tu negocio digital:} Si tu e-commerce no tiene HTTPS, los navegadores muestran ``No seguro'' en rojo. Los clientes no compran. Google te penaliza en el ranking. Es el m\'inimo indispensable.
\end{cajaNota}

% --- Firma digital ---
\subsection{Firma digital: autenticidad + integridad}

\begin{cajaDefinicion}
La \textbf{firma digital} es el equivalente criptogr\'afico de una firma manuscrita, pero m\'as poderosa: garantiza \textbf{autenticidad} (qui\'en firm\'o) e \textbf{integridad} (que el documento no fue modificado despu\'es de firmarse).
\end{cajaDefinicion}

\textbf{¿C\'omo funciona (conceptualmente)?}
\begin{enumerate}[leftmargin=2em]
  \item Se calcula el \textbf{hash} del documento.
  \item Se cifra ese hash con la \textbf{clave privada} del firmante → eso es la firma.
  \item El receptor calcula el hash del documento recibido y lo compara con el hash descifrado (usando la \textbf{clave p\'ublica} del firmante).
  \item Si coinciden → el documento es aut\'entico e \'integro.
\end{enumerate}

\begin{cajaEjemplo}
\textbf{En fintech:} Un contrato de pr\'estamo digital se firma electr\'onicamente. Si alguien modifica una coma del PDF despu\'es de la firma, la verificaci\'on falla y se detecta la alteraci\'on.
\end{cajaEjemplo}

\newpage
% ============================================================
%  SECCIÓN 5 — CASOS DE USO EN NEGOCIOS DIGITALES
% ============================================================
\section{Casos de uso: ¿d\'onde vive la criptograf\'ia en tu negocio?}

\renewcommand{\arraystretch}{1.6}
\begin{tabularx}{\textwidth}{|>{\bfseries}p{3.5cm}|X|>{\centering\arraybackslash}p{2.8cm}|}
\hline
\textbf{Escenario} & \textbf{Criptografía aplicada} & \textbf{Pilar CID} \\
\hline
Cliente paga con tarjeta en un e-commerce & HTTPS/TLS cifra la conexi\'on; la pasarela de pago tokeniza el n\'umero de tarjeta & Confidencialidad \\
\hline
Almacenamiento de contrase\~nas de usuarios & Hash con bcrypt/Argon2 + salt \'unico por usuario & Confidencialidad \\
\hline
Factura electr\'onica o contrato digital & Firma digital con clave privada del emisor & Integridad + Autenticidad \\
\hline
Backup de la base de datos & Cifrado sim\'etrico (AES-256) del archivo de backup & Confidencialidad \\
\hline
Actualizaci\'on de software al cliente & Hash SHA-256 del instalador para verificar que no fue adulterado & Integridad \\
\hline
App de mensajer\'ia (WhatsApp, Signal) & Cifrado de extremo a extremo (E2EE): solo emisor y receptor pueden leer & Confidencialidad \\
\hline
Login con 2FA & C\'odigo TOTP generado con algoritmo HMAC (basado en hash) & Confidencialidad + Disponibilidad \\
\hline
\end{tabularx}

\begin{cajaConcepto}[title={Pregunta para reflexionar}]
Si tu negocio digital no implementa ninguna de estas medidas, ¿qu\'e pilares CID est\'as dejando desprotegidos? ¿Qu\'e impacto tendr\'ia una brecha?
\end{cajaConcepto}

\newpage
% ============================================================
%  SECCIÓN 6 — ACTIVIDAD PRÁCTICA: TALLER DE HASHING
% ============================================================
\section{Actividad pr\'actica: Taller de Hashing}

\begin{cajaActividad}[title={Generar y analizar hashes SHA-256}]

\textbf{Paso 1:} Ingresen a una herramienta online de hashing, por ejemplo:\\
\texttt{https://passwordsgenerator.net/sha256-hash-generator/}\\
(o busquen ``SHA-256 online generator'' en el navegador)

\vspace{6pt}
\textbf{Paso 2:} Generen el hash de las siguientes tres contrase\~nas y completen la tabla:

\vspace{4pt}
\renewcommand{\arraystretch}{1.8}
\begin{tabularx}{\textwidth}{|>{\bfseries}p{4cm}|X|}
\hline
\textbf{Contrase\~na} & \textbf{Hash SHA-256 resultante} \\
\hline
\texttt{123456} & \\[1.5cm]
\hline
\texttt{MiEmpresa2024} & \\[1.5cm]
\hline
\texttt{Xk9\#mL2\$pQ7@vR} & \\[1.5cm]
\hline
\end{tabularx}

\vspace{8pt}
\textbf{Paso 3:} Respondan en grupo:

\vspace{4pt}
\renewcommand{\arraystretch}{1.5}
\begin{tabularx}{\textwidth}{|X|}
\hline
\textbf{¿Qu\'e observan sobre la longitud de los tres hashes? ¿Son todos iguales en largo? ¿Por qu\'e?} \\[2cm]
\hline
\textbf{Si cambio una sola letra de ``MiEmpresa2024'' a ``MiEmpresa2025'', ¿el hash se parece al anterior o es totalmente diferente? ¿Por qu\'e importa esto?} \\[2cm]
\hline
\textbf{¿Pueden ``revertir'' el hash para obtener la contrase\~na original? ¿Qu\'e implica esto para el almacenamiento seguro?} \\[2cm]
\hline
\textbf{¿Por qu\'e la contraseña ``123456'' sigue siendo vulnerable aunque la guarde hasheada?} \newline
\footnotesize{(Pista: piensen en ataques de diccionario y rainbow tables)} \\[2cm]
\hline
\end{tabularx}

\end{cajaActividad}

\newpage
% ============================================================
%  SECCIÓN 7 — ACTIVIDAD CON IA
% ============================================================
\section{Actividad con IA: Explicar, auditar y verificar}

\begin{cajaIA}

\textbf{Objetivo:} Usar una IA generativa como asistente de estudio y luego \textbf{auditar cr\'iticamente} su respuesta. La IA no reemplaza el criterio; lo complementa.

\vspace{8pt}
\textbf{Consigna grupal:}

\begin{enumerate}[leftmargin=2em]
  \item Elijan una IA generativa (ChatGPT, Claude, Gemini, la que prefieran).
  \item P\'idanle exactamente este prompt (o una versi\'on propia):

  \begin{tcolorbox}[colback=white, colframe=violetaIA!50, arc=2mm]
  \small\textit{``Explic\'a en lenguaje simple, como si le hablaras a un estudiante de negocios digitales sin conocimientos t\'ecnicos previos, c\'omo funciona el cifrado de extremo a extremo en una app de mensajería como WhatsApp. Despu\'es, gener\'a un mini-glosario de 8 t\'erminos criptogr\'aficos clave con definiciones de una línea cada una.''}
  \end{tcolorbox}

  \item \textbf{Copien} la respuesta completa de la IA en el espacio de abajo (o impr\'imanla).

  \item \textbf{Auditen} la respuesta contrast\'andola con lo que vieron en clase y con una fuente confiable (material de c\'atedra, ISO/IEC 27000, o una b\'usqueda en sitios como EFF, Cloudflare Learning, etc.).

  \item \textbf{Marquen} con una \ding{51} lo que es correcto, una \ding{55} lo que es impreciso o incorrecto, y una \ding{43} lo que est\'a excesivamente simplificado.

  \item Redacten un \textbf{informe de verificaci\'on de una p\'agina} que incluya:
  \begin{itemize}[leftmargin=1.5em]
    \item ¿Qu\'e afirmaciones de la IA eran correctas?
    \item ¿Qu\'e afirmaciones eran imprecisas o incorrectas?
    \item ¿Qu\'e simplific\'o de más y perdi\'o matiz importante?
    \item ¿Qu\'e corregirían o agregarían?
  \end{itemize}
\end{enumerate}

\end{cajaIA}

\vspace{10pt}

\begin{cajaActividad}[title={Espacio para pegar la respuesta de la IA y auditarla}]

\vspace{4pt}
\renewcommand{\arraystretch}{1.5}
\begin{tabularx}{\textwidth}{|X|}
\hline
\textbf{Pegar aquí la respuesta de la IA (o resumen):} \\[8cm]
\hline
\end{tabularx}

\vspace{6pt}
\begin{tabularx}{\textwidth}{|X|}
\hline
\textbf{Hallazgos de la auditoría (¿qu\'e es correcto? ¿qu\'e es impreciso? ¿qu\'e falta?):} \\[6cm]
\hline
\end{tabularx}

\end{cajaActividad}

\begin{cajaNota}[title={Reglas del curso para el uso de IA}]
\begin{enumerate}[leftmargin=2em]
  \item \textbf{Transparencia:} declarar qu\'e herramienta usaron, para qu\'e y con qu\'e prompt.
  \item \textbf{Verificaci\'on obligatoria:} ninguna salida de IA se entrega ``tal cual''.
  \item \textbf{Responsabilidad:} el grupo es responsable del resultado final, use o no IA.
\end{enumerate}
\end{cajaNota}

\newpage
% ============================================================
%  SECCIÓN 8 — GLOSARIO CRIPTOGRÁFICO
% ============================================================
\section{Glosario criptogr\'afico mínimo}

\renewcommand{\arraystretch}{1.4}
\begin{tabularx}{\textwidth}{>{\bfseries}p{4.5cm}X}
\toprule
\textbf{T\'ermino} & \textbf{Definici\'on en una línea} \\
\midrule
Texto plano & Informaci\'on legible antes de ser cifrada. \\
Texto cifrado & Informaci\'on ilegible resultado del cifrado. \\
Clave sim\'etrica & \'Unica clave compartida para cifrar y descifrar. \\
Clave p\'ublica & Parte del par asim\'etrico que se comparte libremente. \\
Clave privada & Parte del par asim\'etrico que se guarda en secreto. \\
Hash & Funci\'on irreversible que produce una cadena fija desde cualquier dato. \\
SHA-256 & Algoritmo hash que produce una cadena de 256 bits (64 caracteres hex). \\
Sal (salt) & Valor aleatorio agregado antes de hashear para evitar ataques de diccionario. \\
HTTPS / TLS & Protocolo que cifra la comunicaci\'on entre navegador y servidor. \\
Certificado digital & Documento que vincula una clave p\'ublica con una identidad verificada. \\
Firma digital & Hash cifrado con clave privada; garantiza autor\'ia e integridad. \\
Cifrado E2EE & Cifrado de extremo a extremo: solo emisor y receptor pueden leer. \\
Rainbow table & Tabla precomputada de hashes usada para revertir hashes d\'ebiles. \\
bcrypt / Argon2 & Funciones hash lentas diseñadas para almacenar contrase\~nas. \\
\bottomrule
\end{tabularx}

\newpage
% ============================================================
%  SECCIÓN 9 — CIERRE DE UNIDAD 1
% ============================================================
\section{Cierre de Unidad 1: ¿Qu\'e nos llevamos?}

Con esta clase \textbf{cerramos formalmente la Unidad 1}. Hagamos un repaso de todo lo construido:

\begin{cajaCierre}[title={Mapa de la Unidad 1}]

\begin{center}
\begin{tikzpicture}[
  node distance=1.2cm and 2cm,
  box/.style={rectangle, draw=azulUCA, fill=azulUCA!10,
              text width=3.2cm, align=center, rounded corners=3pt,
              font=\small\bfseries},
  boxsmall/.style={rectangle, draw=azulCielo, fill=azulCielo!10,
              text width=2.8cm, align=center, rounded corners=3pt,
              font=\scriptsize},
  arr/.style={-{Stealth[length=2.5mm]}, thick, azulUCA}
]
  % Fila 1: Fundamentos
  \node[box] (cid) {PILARES CID\\Confidencialidad\\Integridad\\Disponibilidad};
  \node[box, right=2.5cm of cid] (conc) {CONCEPTOS\\Activo, Amenaza,\\Vulnerabilidad,\\Impacto, Riesgo};
  \node[box, right=2.5cm of conc] (amen) {AMENAZAS\\Phishing, Ransomware,\\Ing. Social, Fraude};

  % Fila 2: Herramientas y cripto
  \node[boxsmall, below=1.5cm of cid] (herr) {HERRAMIENTAS\\HIBP, VirusTotal,\\Gestores de contrase\~nas};
  \node[boxsmall, below=1.5cm of conc] (cripto) {CRIPTOGRAFÍA\\Cifrado, Hash,\\HTTPS, Firma digital};
  \node[boxsmall, below=1.5cm of amen] (caso) {CASOS\\E-commerce,\\Fintech, Apps};

  % Flechas
  \draw[arr] (cid.south) -- (herr.north);
  \draw[arr] (conc.south) -- (cripto.north);
  \draw[arr] (amen.south) -- (caso.north);
  \draw[arr] (cid.east) -- (conc.west);
  \draw[arr] (conc.east) -- (amen.west);
  \draw[arr, dashed] (herr.east) -- (cripto.west);
  \draw[arr, dashed] (cripto.east) -- (caso.west);
\end{tikzpicture}
\end{center}

\vspace{8pt}
\textbf{En la Unidad 1 aprendimos a:}
\begin{itemize}[leftmargin=2em]
  \item Identificar activos y clasificar riesgos usando los pilares CID.
  \item Reconocer las amenazas m\'as comunes para negocios digitales.
  \item Usar herramientas gratuitas de verificaci\'on y protecci\'on.
  \item Comprender c\'omo la criptograf\'ia habilita la confianza digital.
  \item Usar IA generativa de forma \textbf{cr\'itica y verificada}.
\end{itemize}

\vspace{6pt}
\textbf{Lo que viene (Unidad 2):} Gesti\'on de riesgos formal. Vamos a pasar de ``entender las amenazas'' a \textbf{medirlas, priorizarlas y tratarlas} con metodología (matriz P$\times$I, BIA, controles, ISO 27001).
\end{cajaCierre}

\newpage
% ============================================================
%  SECCIÓN 10 — NOTAS
% ============================================================
\section{Espacio para notas y reflexiones}

\begin{tcolorbox}[colback=grisClaro, colframe=gray!50, arc=2mm, height=7cm]
\textcolor{gray}{\itshape Us\'a este espacio para anotar dudas, ideas para tu TP Integrador, o conexiones entre la criptograf\'ia y el negocio digital que eligi\'o tu grupo.}
\end{tcolorbox}

\vspace{12pt}

\begin{cajaConcepto}[title={Puente hacia la Clase 5}]
La pr\'oxima clase arrancamos la \textbf{Unidad 2: Gesti\'on de Riesgos}.\\[4pt]
Pregunta disparadora para pensar antes de la pr\'oxima clase:\\
\emph{Si tuvieras que proteger tu negocio digital con presupuesto para solo 3 controles de seguridad, ¿cu\'ales elegir\'ias y por qu\'e?}
\end{cajaConcepto}

\vfill
\begin{center}
\small\textcolor{gray}{%
Cuadernillo elaborado para la c\'atedra de Seguridad de la Informaci\'on y de
Sistemas --- Lic.\ en Gesti\'on de Negocios Digitales, UCA.\\
Prohibida su distribuci\'on fuera del \'ambito de la cursada sin autorizaci\'on del docente.}
\end{center}

\end{document}